% Fichier généré par build_notebook_appendix.py — ne pas éditer à la main.
\subsection{\texttt{src/nb\_common.py}}
\begin{lstlisting}[style=py]
"""
nb_common.py
============

Briques *communes* aux deux implémentations de Naive Bayes multinomial
(version RDD dans ``nb_rdd.py`` et version DataFrames dans ``nb_dataframe.py``).

On centralise ici tout ce qui **doit être strictement identique** entre les deux
versions pour garantir des prédictions rigoureusement égales :

1. la tokenisation (bag-of-words) ;
2. la construction du vocabulaire et de la table des étiquettes ;
3. la vectorisation d'un document en liste d'indices de mots ;
4. la *construction du modèle* à partir de comptes entiers (comptes -> log-probas) ;
5. la *fonction de prédiction* d'un document à partir du modèle.

Les points 4 et 5 sont volontairement écrits **une seule fois** : les versions RDD
et DataFrames se contentent de produire, chacune à leur manière (map/reduce),
exactement les mêmes comptes entiers, puis appellent ces fonctions communes.
Comme les comptes sont des entiers (donc reproductibles au bit près) et que la
prédiction est déterministe, les deux versions renvoient des prédictions
identiques.

La modélisation reproduit fidèlement ``sklearn.naive_bayes.MultinomialNB`` avec
``alpha=1`` et ``fit_prior=True`` :

    log P(c)      = log( N_c / N )
    log P(w | c)  = log( (count_{w,c} + alpha) / (total_c + alpha * V) )

où :
    N            = nombre de documents d'entraînement
    N_c          = nombre de documents de la classe c
    count_{w,c}  = nombre total d'occurrences du mot w dans les docs de classe c
    total_c      = sum_w count_{w,c}  (nombre total de tokens de la classe c)
    V            = taille du vocabulaire
    alpha        = paramètre de lissage de Laplace (1 par défaut)

Score d'un document d pour une classe c (tout en logarithme) :

    score(d, c) = log P(c) + sum_{w in d} count_w(d) * log P(w | c)

La classe prédite est celle qui maximise ce score.
"""

from __future__ import annotations

import math
import os
import re
import subprocess
from dataclasses import dataclass, field
from typing import Dict, List, Sequence, Tuple

# ---------------------------------------------------------------------------
# 1. Tokenisation (bag-of-words)
# ---------------------------------------------------------------------------
# On reproduit la tokenisation par défaut de sklearn.CountVectorizer :
#   - passage en minuscules,
#   - motif r"\b\w\w+\b" : suites d'au moins 2 caractères de mot.
# Utiliser exactement la même règle est indispensable pour pouvoir comparer nos
# résultats à ceux de sklearn.MultinomialNB dans le notebook.
_TOKEN_RE = re.compile(r"(?u)\b\w\w+\b")


def tokenize(text: str) -> List[str]:
    """Découpe un texte en liste de tokens (mots), comme sklearn CountVectorizer.

    >>> tokenize("Free entry, WIN a prize!!!")
    ['free', 'entry', 'win', 'prize']
    """
    if text is None:
        return []
    return _TOKEN_RE.findall(text.lower())


# ---------------------------------------------------------------------------
# 2. Vocabulaire et étiquettes
# ---------------------------------------------------------------------------
def build_vocabulary(
    tokenized_docs: Sequence[Sequence[str]], min_df: int = 1
) -> Dict[str, int]:
    """Construit le vocabulaire {mot -> indice} à partir des documents d'entraînement.

    Le vocabulaire est trié par ordre alphabétique (comme sklearn) afin que les
    indices soient parfaitement déterministes d'une exécution à l'autre.

    Paramètres
    ----------
    tokenized_docs : documents déjà tokenisés (liste de listes de mots).
    min_df         : nombre minimal de documents dans lesquels un mot doit
                     apparaître pour être conservé (1 = on garde tout).
    """
    # Compte le nombre de *documents* contenant chaque mot (document frequency).
    doc_freq: Dict[str, int] = {}
    for tokens in tokenized_docs:
        for w in set(tokens):  # set() -> un mot compte une seule fois par doc
            doc_freq[w] = doc_freq.get(w, 0) + 1

    # On ne garde que les mots suffisamment fréquents, puis on trie -> indices stables.
    kept = sorted(w for w, df in doc_freq.items() if df >= min_df)
    return {w: i for i, w in enumerate(kept)}


def build_label_index(labels: Sequence[str]) -> Dict[str, int]:
    """Construit la table {étiquette -> entier} (triée pour être déterministe)."""
    return {lab: i for i, lab in enumerate(sorted(set(labels)))}


def doc_to_indices(tokens: Sequence[str], vocab: Dict[str, int]) -> List[int]:
    """Vectorise un document en liste d'indices de mots (avec répétitions).

    Les mots hors-vocabulaire (OOV) sont ignorés : cela correspond exactement au
    comportement de sklearn, où la matrice de features a un nombre fixe de
    colonnes défini au moment du ``fit`` sur l'ensemble d'entraînement.

    Exemple : si vocab = {"free": 0, "win": 1} et tokens = ["free", "win", "free", "zzz"],
    on renvoie [0, 1, 0] (le token OOV "zzz" est retiré).
    """
    out: List[int] = []
    for w in tokens:
        idx = vocab.get(w)
        if idx is not None:
            out.append(idx)
    return out


# ---------------------------------------------------------------------------
# 3. Le modèle Naive Bayes (structure partagée)
# ---------------------------------------------------------------------------
@dataclass
class NaiveBayesModel:
    """Modèle Naive Bayes multinomial *dense-équivalent mais stocké creux*.

    Pour respecter le lissage de Laplace, sklearn stocke une log-proba pour
    CHAQUE (classe, mot) du vocabulaire, y compris les mots de compte nul dans
    une classe. Ce serait coûteux (n_classes * V valeurs). On stocke donc :

      - ``log_prior[c]``              : log P(c)
      - ``log_likelihood[c][w]``      : log P(w|c) uniquement pour les mots dont
                                        le compte est NON nul dans la classe c
      - ``log_likelihood_default[c]`` : log P(w|c) pour tout mot de compte nul
                                        dans la classe c
                                        = log( alpha / (total_c + alpha * V) )

    C'est rigoureusement équivalent au calcul dense de sklearn : lors de la
    prédiction, un mot présent dans le document mais de compte nul dans la
    classe utilise simplement la valeur ``default`` (voir ``predict_indices``).
    Cette représentation creuse est aussi ce que l'on *broadcast* aux workers.
    """

    n_classes: int
    n_docs: int
    vocab_size: int
    alpha: float
    # index entier de classe -> étiquette d'origine (str)
    idx_to_label: List[str]
    log_prior: List[float]
    # log_likelihood[c] : dict {indice_mot -> log P(w|c)} (comptes non nuls seult.)
    log_likelihood: List[Dict[int, float]]
    log_likelihood_default: List[float]


def build_model(
    *,
    n_docs: int,
    vocab_size: int,
    idx_to_label: List[str],
    class_doc_counts: Dict[int, int],
    class_token_totals: Dict[int, int],
    word_class_counts: Dict[Tuple[int, int], int],
    alpha: float = 1.0,
) -> NaiveBayesModel:
    """Construit le modèle (log-probas) à partir de **comptes entiers**.

    C'est le cœur du calcul Naive Bayes, factorisé ici pour que les versions RDD
    et DataFrames produisent EXACTEMENT le même modèle : il leur suffit de fournir
    les mêmes agrégats entiers (calculés en map/reduce), le reste est identique.

    Paramètres (tous issus d'un comptage MapReduce dans les versions Spark)
    ----------------------------------------------------------------------
    n_docs             : N   (nb total de documents d'entraînement)
    vocab_size         : V   (taille du vocabulaire)
    idx_to_label       : correspondance indice_classe -> étiquette
    class_doc_counts   : {c -> N_c}         nb de docs par classe
    class_token_totals : {c -> total_c}     nb total de tokens par classe
    word_class_counts  : {(c, w) -> count}  nb d'occurrences du mot w en classe c
    alpha              : lissage de Laplace (1.0 par défaut)
    """
    n_classes = len(idx_to_label)

    # --- Priors : log P(c) = log(N_c / N) -----------------------------------
    log_prior = [0.0] * n_classes
    for c in range(n_classes):
        n_c = class_doc_counts.get(c, 0)
        # Un prior nul (classe absente) donnerait log(0) = -inf ; en pratique
        # toutes les classes sont présentes en entraînement.
        log_prior[c] = math.log(n_c / n_docs) if n_c > 0 else float("-inf")

    # --- Log-vraisemblance par mot : log P(w|c) -----------------------------
    # Dénominateur commun à une classe : total_c + alpha * V
    denom = [class_token_totals.get(c, 0) + alpha * vocab_size for c in range(n_classes)]

    # Valeur par défaut (mot de compte nul dans la classe) : log(alpha / denom)
    log_likelihood_default = [
        math.log(alpha / denom[c]) if denom[c] > 0 else float("-inf")
        for c in range(n_classes)
    ]

    # Pour chaque (classe, mot) réellement observé : log((count + alpha) / denom)
    log_likelihood: List[Dict[int, float]] = [dict() for _ in range(n_classes)]
    for (c, w), count in word_class_counts.items():
        log_likelihood[c][w] = math.log((count + alpha) / denom[c])

    return NaiveBayesModel(
        n_classes=n_classes,
        n_docs=n_docs,
        vocab_size=vocab_size,
        alpha=alpha,
        idx_to_label=list(idx_to_label),
        log_prior=log_prior,
        log_likelihood=log_likelihood,
        log_likelihood_default=log_likelihood_default,
    )


def predict_indices(model: NaiveBayesModel, token_indices: Sequence[int]) -> int:
    """Prédit l'indice de classe d'un document (liste d'indices de mots).

    Fonction PARTAGÉE par les deux versions -> prédictions identiques garanties.

    On calcule, pour chaque classe c :

        score(c) = log P(c) + sum_{w in doc} log P(w|c)

    (chaque occurrence d'un mot ajoute son log P(w|c), ce qui revient bien à
    multiplier par le compte du mot dans le document). On renvoie l'argmax.
    """
    best_c = 0
    best_score = float("-inf")
    for c in range(model.n_classes):
        score = model.log_prior[c]
        ll_c = model.log_likelihood[c]
        default_c = model.log_likelihood_default[c]
        for w in token_indices:
            # get(w, default) : mot de compte nul dans la classe -> valeur lissée.
            score += ll_c.get(w, default_c)
        if score > best_score:
            best_score = score
            best_c = c
    return best_c


def accuracy(y_true: Sequence[int], y_pred: Sequence[int]) -> float:
    """Exactitude (fraction de prédictions correctes)."""
    if not y_true:
        return 0.0
    correct = sum(1 for a, b in zip(y_true, y_pred) if a == b)
    return correct / len(y_true)


# ---------------------------------------------------------------------------
# 4. Chargement des jeux de données
# ---------------------------------------------------------------------------
def load_sms_spam(csv_path: str) -> Tuple[List[str], List[str]]:
    """Charge le jeu SMS Spam Collection depuis un CSV (colonnes label,text).

    Le fichier est produit par ``data/download_sms_spam.py``.
    Renvoie (textes, étiquettes) avec étiquettes dans {"ham", "spam"}.
    """
    import csv

    texts: List[str] = []
    labels: List[str] = []
    with open(csv_path, "r", encoding="utf-8", newline="") as fh:
        reader = csv.DictReader(fh)
        for row in reader:
            labels.append(row["label"])
            texts.append(row["text"])
    return texts, labels


def load_20newsgroups(
    categories: List[str] | None = None, subset: str = "all"
) -> Tuple[List[str], List[str]]:
    """Charge 20 Newsgroups via scikit-learn (téléchargé/caché automatiquement).

    Renvoie (textes, étiquettes) où l'étiquette est le nom de la catégorie.
    """
    from sklearn.datasets import fetch_20newsgroups

    bunch = fetch_20newsgroups(
        subset=subset,
        categories=categories,
        remove=("headers", "footers", "quotes"),  # pour un signal purement textuel
        shuffle=True,
        random_state=42,
    )
    labels = [bunch.target_names[t] for t in bunch.target]
    return list(bunch.data), labels


def replicate_dataset(
    texts: List[str], labels: List[str], factor: int
) -> Tuple[List[str], List[str]]:
    """Duplique le jeu de données ``factor`` fois (pour l'étude de scalabilité).

    Multiplier la taille des données à distribution constante permet de mesurer
    comment le temps d'exécution évolue avec le *volume*, sans changer l'accuracy.
    """
    if factor <= 1:
        return list(texts), list(labels)
    return texts * factor, labels * factor


# ---------------------------------------------------------------------------
# 5. Découpage entraînement / test
# ---------------------------------------------------------------------------
def train_test_split_texts(
    texts: List[str],
    labels: List[str],
    test_size: float = 0.2,
    seed: int = 42,
) -> Tuple[List[str], List[str], List[str], List[str]]:
    """Découpe (textes, étiquettes) en ensembles d'entraînement et de test.

    On délègue à sklearn pour un découpage stratifié et reproductible.
    Renvoie (X_train, X_test, y_train, y_test).
    """
    from sklearn.model_selection import train_test_split

    # stratify=labels garde les mêmes proportions de classes dans train et test.
    X_train, X_test, y_train, y_test = train_test_split(
        texts, labels, test_size=test_size, random_state=seed, stratify=labels
    )
    return X_train, X_test, y_train, y_test


# ---------------------------------------------------------------------------
# 6. Préparation commune : textes -> (indices de classe, indices de mots)
# ---------------------------------------------------------------------------
@dataclass
class PreparedData:
    """Données prêtes pour l'entraînement Spark, partagées par les deux versions."""

    vocab: Dict[str, int]
    label_to_idx: Dict[str, int]
    idx_to_label: List[str]
    vocab_size: int
    # Chaque exemple : (indice_classe, [indices de mots avec répétitions])
    train: List[Tuple[int, List[int]]]
    test: List[Tuple[int, List[int]]]


def prepare(
    X_train: List[str],
    y_train: List[str],
    X_test: List[str],
    y_test: List[str],
    min_df: int = 1,
) -> PreparedData:
    """Tokenise, construit vocabulaire + étiquettes, et vectorise train/test.

    Cette étape déterministe est faite dans le driver (les données tiennent en
    mémoire ; c'est du prétraitement). Le vocabulaire est construit **uniquement
    sur l'entraînement** (comme un ``fit`` sklearn), puis appliqué au test.
    """
    train_tokens = [tokenize(t) for t in X_train]
    test_tokens = [tokenize(t) for t in X_test]

    vocab = build_vocabulary(train_tokens, min_df=min_df)
    label_to_idx = build_label_index(y_train)
    idx_to_label = [lab for lab, _ in sorted(label_to_idx.items(), key=lambda kv: kv[1])]

    train = [
        (label_to_idx[lab], doc_to_indices(toks, vocab))
        for lab, toks in zip(y_train, train_tokens)
    ]
    # Les étiquettes de test inconnues à l'entraînement sont improbables ici
    # (split stratifié) ; on suppose qu'elles existent dans label_to_idx.
    test = [
        (label_to_idx[lab], doc_to_indices(toks, vocab))
        for lab, toks in zip(y_test, test_tokens)
    ]

    return PreparedData(
        vocab=vocab,
        label_to_idx=label_to_idx,
        idx_to_label=idx_to_label,
        vocab_size=len(vocab),
        train=train,
        test=test,
    )


# ---------------------------------------------------------------------------
# 7. SparkSession avec le bon JDK
# ---------------------------------------------------------------------------
def _resolve_java_home() -> str | None:
    """Trouve un JAVA_HOME compatible Spark 3.5 (JDK 8/11/17).

    Spark 3.5 ne supporte PAS les JDK plus récents (20, 21, 23...). Sur cette
    machine le ``java`` par défaut peut être trop récent : on force donc un JDK 17
    via ``/usr/libexec/java_home -v 17`` (macOS).
    """
    # Si JAVA_HOME est déjà positionné sur un JDK supporté, on le respecte.
    for version in ("17", "11", "1.8"):
        try:
            out = subprocess.run(
                ["/usr/libexec/java_home", "-v", version],
                capture_output=True,
                text=True,
                check=True,
            )
            path = out.stdout.strip()
            if path and os.path.isdir(path):
                return path
        except Exception:
            continue
    return None


def get_spark(app_name: str = "NaiveBayesMapReduce", master: str = "local[*]"):
    """Crée (ou récupère) une SparkSession configurée pour ce projet.

    - force un JDK compatible (17) si nécessaire ;
    - ``master`` permet de choisir le parallélisme (ex. "local[2]", "local[*]").
    """
    java_home = _resolve_java_home()
    if java_home:
        os.environ["JAVA_HOME"] = java_home

    # Le worker Python doit utiliser le même interpréteur que le driver (venv).
    import sys

    os.environ.setdefault("PYSPARK_PYTHON", sys.executable)
    os.environ.setdefault("PYSPARK_DRIVER_PYTHON", sys.executable)

    from pyspark.sql import SparkSession

    spark = (
        SparkSession.builder.appName(app_name)
        .master(master)
        # Logs moins verbeux et démarrage plus rapide en local.
        .config("spark.ui.showConsoleProgress", "false")
        .config("spark.sql.shuffle.partitions", "8")
        .getOrCreate()
    )

    # Les workers exécutent nos fonctions (predict_indices) et désérialisent le
    # modèle (dataclass NaiveBayesModel) : ils doivent donc pouvoir importer nos
    # modules. On expédie les fichiers source du projet à tous les exécuteurs.
    src_dir = os.path.dirname(os.path.abspath(__file__))
    for mod in ("nb_common.py", "nb_rdd.py", "nb_dataframe.py"):
        path = os.path.join(src_dir, mod)
        if os.path.isfile(path):
            spark.sparkContext.addPyFile(path)

    return spark
\end{lstlisting}
\subsection{\texttt{src/nb\_rdd.py}}
\begin{lstlisting}[style=py]
"""
nb_rdd.py
=========

Naive Bayes multinomial **en MapReduce sur des RDD Spark**.

C'est la version « bas niveau » : on manipule directement des RDD et on exprime
l'entraînement comme une suite d'opérations *map* / *reduceByKey*, qui est
exactement l'esprit MapReduce.

Idée clé — comptages par MapReduce :
    Un document est (indice_classe, [indices de mots]).
    * MAP    : pour chaque occurrence d'un mot w dans un doc de classe c,
               on émet la paire clé/valeur ((c, w), 1).
    * REDUCE : reduceByKey(add) additionne les 1 -> nombre d'occurrences de w
               dans la classe c, soit count_{w,c}.
    De même on compte les documents par classe et les tokens par classe.

Ces agrégats entiers sont ensuite passés à ``nb_common.build_model`` (calcul des
log-probas), puis le modèle est **broadcasté** aux workers pour la prédiction.
"""

from __future__ import annotations

from operator import add
from typing import List, Tuple

import nb_common as C


def train_rdd(sc, train_data: List[Tuple[int, List[int]]], vocab_size: int,
              idx_to_label: List[str], alpha: float = 1.0, num_partitions: int = 8
              ) -> C.NaiveBayesModel:
    """Entraîne le modèle Naive Bayes en RDD (map/reduceByKey).

    Paramètres
    ----------
    sc             : SparkContext.
    train_data     : liste de (indice_classe, [indices de mots]).
    vocab_size     : V, taille du vocabulaire.
    idx_to_label   : correspondance indice_classe -> étiquette d'origine.
    alpha          : lissage de Laplace.
    num_partitions : nombre de partitions du RDD (parallélisme des données).
    """
    # On distribue les documents sur le cluster (ici local[k]) en `num_partitions`.
    docs = sc.parallelize(train_data, numSlices=num_partitions)

    # --- (a) Nombre total de documents : N ----------------------------------
    n_docs = docs.count()

    # --- (b) Nombre de documents par classe : N_c ---------------------------
    #   MAP    : chaque doc -> (classe, 1)
    #   REDUCE : reduceByKey(add) -> somme des 1 par classe
    class_doc_counts = dict(
        docs.map(lambda cw: (cw[0], 1)).reduceByKey(add).collect()
    )

    # --- (c) Nombre total de tokens par classe : total_c --------------------
    #   MAP    : chaque doc -> (classe, nb de tokens du doc)
    #   REDUCE : somme par classe
    class_token_totals = dict(
        docs.map(lambda cw: (cw[0], len(cw[1]))).reduceByKey(add).collect()
    )

    # --- (d) Compte des mots par classe : count_{w,c}  (LE cœur MapReduce) ---
    #   MAP (flatMap) : pour un doc (c, [w1, w2, w1, ...]), on émet
    #                   ((c, w1), 1), ((c, w2), 1), ((c, w1), 1), ...
    #                   -> une paire par occurrence de mot.
    #   REDUCE        : reduceByKey(add) additionne toutes les occurrences,
    #                   donnant count_{w,c} pour chaque couple (classe, mot).
    def emit_word_class(doc: Tuple[int, List[int]]):
        c, word_indices = doc
        for w in word_indices:
            yield ((c, w), 1)

    word_class_counts = dict(
        docs.flatMap(emit_word_class).reduceByKey(add).collect()
    )

    # --- (e) Construction du modèle (comptes entiers -> log-probas) ---------
    # Calcul factorisé dans nb_common pour être identique à la version DataFrame.
    return C.build_model(
        n_docs=n_docs,
        vocab_size=vocab_size,
        idx_to_label=idx_to_label,
        class_doc_counts=class_doc_counts,
        class_token_totals=class_token_totals,
        word_class_counts=word_class_counts,
        alpha=alpha,
    )


def predict_rdd(sc, model: C.NaiveBayesModel,
                test_data: List[Tuple[int, List[int]]], num_partitions: int = 8
                ) -> List[int]:
    """Prédit les classes des documents de test, en parallèle sur les RDD.

    Le modèle est **broadcasté** : Spark l'envoie une seule fois à chaque worker
    (au lieu de le sérialiser avec chaque tâche), ce qui est essentiel quand le
    modèle est volumineux (V mots x C classes).
    """
    # BROADCAST : diffusion du modèle en lecture seule à tous les workers.
    bc_model = sc.broadcast(model)

    docs = sc.parallelize(test_data, numSlices=num_partitions)

    # MAP : chaque doc -> classe prédite, via la fonction de prédiction COMMUNE.
    # bc_model.value récupère le modèle local au worker (pas de re-sérialisation).
    preds = docs.map(
        lambda cw: C.predict_indices(bc_model.value, cw[1])
    ).collect()

    bc_model.unpersist()
    return preds


def evaluate_rdd(sc, model: C.NaiveBayesModel,
                 test_data: List[Tuple[int, List[int]]], num_partitions: int = 8
                 ) -> float:
    """Renvoie l'exactitude (accuracy) du modèle sur l'ensemble de test."""
    y_true = [c for c, _ in test_data]
    y_pred = predict_rdd(sc, model, test_data, num_partitions=num_partitions)
    return C.accuracy(y_true, y_pred)
\end{lstlisting}
\subsection{\texttt{src/nb\_dataframe.py}}
\begin{lstlisting}[style=py]
"""
nb_dataframe.py
===============

Naive Bayes multinomial **avec l'API DataFrames de Spark**.

Même algorithme que ``nb_rdd.py``, mais exprimé avec des opérations
relationnelles de haut niveau (``explode``, ``groupBy`` / ``agg``) au lieu de
map/reduceByKey. Le moteur Catalyst optimise ces opérations.

Correspondance avec la version RDD :
    * ``explode`` d'un tableau de mots  ≈  la phase MAP (une ligne par mot) ;
    * ``groupBy(...).agg(count/sum)``   ≈  la phase REDUCE (reduceByKey).

On produit exactement les mêmes agrégats entiers que la version RDD, puis on
appelle le MÊME ``nb_common.build_model`` et la MÊME fonction de prédiction :
les deux versions renvoient donc des prédictions identiques.
Le modèle est **broadcasté** aux workers pour la prédiction via une UDF.
"""

from __future__ import annotations

from typing import List, Tuple

from pyspark.sql import Row
from pyspark.sql import functions as F
from pyspark.sql.types import (
    ArrayType,
    IntegerType,
    StructField,
    StructType,
)

import nb_common as C


def _to_dataframe(spark, data: List[Tuple[int, List[int]]], num_partitions: int):
    """Construit un DataFrame (label:int, words:array<int>) à partir des données.

    Chaque ligne = un document : son indice de classe et son sac de mots (liste
    d'indices de mots, avec répétitions).
    """
    schema = StructType([
        StructField("label", IntegerType(), False),
        StructField("words", ArrayType(IntegerType()), False),
    ])
    rows = [Row(label=c, words=list(w)) for c, w in data]
    return spark.createDataFrame(rows, schema=schema).repartition(num_partitions)


def train_dataframe(spark, train_data: List[Tuple[int, List[int]]], vocab_size: int,
                    idx_to_label: List[str], alpha: float = 1.0, num_partitions: int = 8
                    ) -> C.NaiveBayesModel:
    """Entraîne le modèle Naive Bayes avec l'API DataFrames (groupBy/agg)."""
    df = _to_dataframe(spark, train_data, num_partitions)
    df.cache()  # réutilisé plusieurs fois ci-dessous -> on le met en cache

    # --- (a) Nombre total de documents : N ----------------------------------
    n_docs = df.count()

    # --- (b) Nombre de documents par classe : N_c ---------------------------
    #   groupBy(label).count()  ≈  REDUCE : compte les lignes par classe.
    class_doc_counts = {
        r["label"]: r["cnt"]
        for r in df.groupBy("label").agg(F.count("*").alias("cnt")).collect()
    }

    # --- (c) Nombre total de tokens par classe : total_c --------------------
    #   size(words) = nb de tokens du doc ; on somme par classe.
    class_token_totals = {
        r["label"]: r["tot"]
        for r in df.groupBy("label")
        .agg(F.sum(F.size("words")).alias("tot"))
        .collect()
    }

    # --- (d) Compte des mots par classe : count_{w,c}  (cœur MAP/REDUCE) -----
    #   MAP    : explode(words) -> une ligne (label, word) PAR occurrence de mot
    #            (équivaut au flatMap qui émet ((c, w), 1) de la version RDD).
    #   REDUCE : groupBy(label, word).count() -> count_{w,c}.
    exploded = df.select("label", F.explode("words").alias("word"))
    word_class_counts = {
        (r["label"], r["word"]): r["cnt"]
        for r in exploded.groupBy("label", "word")
        .agg(F.count("*").alias("cnt"))
        .collect()
    }

    df.unpersist()

    # --- (e) Construction du modèle (identique à la version RDD) ------------
    return C.build_model(
        n_docs=n_docs,
        vocab_size=vocab_size,
        idx_to_label=idx_to_label,
        class_doc_counts=class_doc_counts,
        class_token_totals=class_token_totals,
        word_class_counts=word_class_counts,
        alpha=alpha,
    )


def predict_dataframe(spark, model: C.NaiveBayesModel,
                      test_data: List[Tuple[int, List[int]]], num_partitions: int = 8
                      ) -> List[int]:
    """Prédit les classes de test via une UDF s'appuyant sur le modèle broadcasté."""
    sc = spark.sparkContext

    # BROADCAST : le modèle est diffusé une fois par worker (lecture seule),
    # au lieu d'être capturé/sérialisé pour chaque tâche.
    bc_model = sc.broadcast(model)

    # UDF : applique la fonction de prédiction COMMUNE à chaque sac de mots.
    # La closure référence bc_model.value, résolu localement sur chaque worker.
    @F.udf(returnType=IntegerType())
    def predict_udf(words):
        return int(C.predict_indices(bc_model.value, words or []))

    # On attache un identifiant EXPLICITE (position dans test_data) AVANT tout
    # shuffle, afin de pouvoir réordonner les prédictions et les aligner sur
    # test_data, quel que soit le repartitionnement effectué par Spark.
    schema = StructType([
        StructField("idx", IntegerType(), False),
        StructField("words", ArrayType(IntegerType()), False),
    ])
    rows_in = [Row(idx=i, words=list(w)) for i, (_, w) in enumerate(test_data)]
    df = spark.createDataFrame(rows_in, schema=schema).repartition(num_partitions)

    df = df.withColumn("pred", predict_udf(F.col("words")))
    rows = df.select("idx", "pred").orderBy("idx").collect()

    bc_model.unpersist()
    return [r["pred"] for r in rows]


def evaluate_dataframe(spark, model: C.NaiveBayesModel,
                       test_data: List[Tuple[int, List[int]]], num_partitions: int = 8
                       ) -> float:
    """Renvoie l'exactitude (accuracy) du modèle sur l'ensemble de test."""
    y_true = [c for c, _ in test_data]
    y_pred = predict_dataframe(spark, model, test_data, num_partitions=num_partitions)
    return C.accuracy(y_true, y_pred)
\end{lstlisting}
\subsection{\texttt{src/benchmark.py}}
\begin{lstlisting}[style=py]
"""
benchmark.py
============

Étude de **scalabilité** de Naive Bayes multinomial en Spark.

On fait varier deux facteurs et on mesure les temps d'entraînement et de
prédiction pour les deux versions (RDD et DataFrames) :

  1. la **taille des données**  : on réplique le jeu 20 Newsgroups plusieurs fois
     (facteurs de réplication), pour observer le passage à l'échelle en volume ;
  2. le **parallélisme**        : on fait varier ``local[k]`` (nombre de cœurs),
     pour observer l'accélération (speed-up) à volume constant.

Deux précautions méthodologiques importantes :

  * **Pas de fuite de données** : la réplication est faite APRÈS le découpage
    train/test, et séparément sur chaque côté. Ainsi une copie d'un document
    d'entraînement ne peut jamais se retrouver dans le test (et inversement).
    L'accuracy reste donc honnête et stable quel que soit le facteur de
    réplication (elle ne doit PAS augmenter avec le volume).
  * **Bruit de mesure** : avec ``--repeat N`` (N > 1), chaque point est mesuré
    N fois ; on reporte la **moyenne** et l'**écart-type** des temps. Cela lisse
    les exécutions aberrantes (variabilité système, GC de la JVM, etc.).

Sorties :
  - ``results/results.csv``               : toutes les mesures (moyennes + écarts-types) ;
  - ``results/scalability_datasize.png``  : temps vs taille des données ;
  - ``results/scalability_cores.png``     : temps vs nombre de cœurs (local[k]).

Usage :
    python src/benchmark.py                                   # configuration par défaut
    python src/benchmark.py --quick                           # version rapide (petits facteurs)
    python src/benchmark.py --factors 1 2 4 8 --cores 1 2 4 8 # benchmark complet
    python src/benchmark.py --factors 1 2 4 8 --cores 1 2 4 8 --repeat 3  # moyenné
"""

from __future__ import annotations

import argparse
import csv
import os
import statistics
import time
from typing import Callable, List, Tuple

import nb_common as C
import nb_rdd
import nb_dataframe

HERE = os.path.dirname(os.path.abspath(__file__))
ROOT = os.path.dirname(HERE)
RESULTS_DIR = os.path.join(ROOT, "results")


def _timeit(fn: Callable):
    """Exécute ``fn`` et renvoie (résultat, temps écoulé en secondes)."""
    t0 = time.perf_counter()
    result = fn()
    return result, time.perf_counter() - t0


def _mean_std(values: List[float]) -> Tuple[float, float]:
    """Renvoie (moyenne, écart-type de population) d'une liste de mesures.

    L'écart-type vaut 0 s'il n'y a qu'une seule mesure (``--repeat 1``).
    """
    mean = statistics.mean(values)
    std = statistics.pstdev(values) if len(values) > 1 else 0.0
    return mean, std


def run_once(master: str, factor: int, base_texts: List[str], base_labels: List[str],
             repeat: int = 1, seed: int = 42) -> List[dict]:
    """Mesure un point de benchmark pour un (master, facteur de réplication) donné.

    Méthodologie anti-fuite : on découpe D'ABORD le jeu de base en train/test, PUIS
    on réplique chaque côté séparément. La réplication grossit donc le volume
    (train et test) sans jamais mélanger un document entre les deux ensembles.

    Avec ``repeat > 1``, on ré-exécute l'entraînement et la prédiction ``repeat``
    fois sur les MÊMES données (la SparkSession est réutilisée), et on agrège les
    temps en moyenne ± écart-type.

    Renvoie une liste de dicts (une ligne par version RDD/DataFrame).
    """
    # 1) Découpage AVANT réplication -> pas de fuite de données.
    X_tr0, X_te0, y_tr0, y_te0 = C.train_test_split_texts(base_texts, base_labels, seed=seed)
    # 2) Réplication INDÉPENDANTE de chaque côté pour faire grossir le volume.
    X_tr, y_tr = C.replicate_dataset(X_tr0, y_tr0, factor)
    X_te, y_te = C.replicate_dataset(X_te0, y_te0, factor)
    data = C.prepare(X_tr, y_tr, X_te, y_te)

    # 3) SparkSession avec le parallélisme voulu (local[k]).
    spark = C.get_spark(app_name=f"bench-{master}-x{factor}", master=master)
    spark.sparkContext.setLogLevel("ERROR")
    sc = spark.sparkContext
    # Nombre de partitions ~ proportionnel au parallélisme (au moins 4).
    parts = max(4, sc.defaultParallelism)

    rows: List[dict] = []
    try:
        n_train = len(data.train)
        n_test = len(data.test)

        # Accumulateurs de temps sur les ``repeat`` exécutions.
        t_train_rdd_runs: List[float] = []
        t_pred_rdd_runs: List[float] = []
        t_train_df_runs: List[float] = []
        t_pred_df_runs: List[float] = []
        acc_rdd = acc_df = None

        for _ in range(repeat):
            # --- Version RDD ------------------------------------------------
            model_rdd, t_train_rdd = _timeit(
                lambda: nb_rdd.train_rdd(sc, data.train, data.vocab_size,
                                         data.idx_to_label, num_partitions=parts)
            )
            acc_rdd, t_pred_rdd = _timeit(
                lambda: nb_rdd.evaluate_rdd(sc, model_rdd, data.test, num_partitions=parts)
            )

            # --- Version DataFrame -----------------------------------------
            model_df, t_train_df = _timeit(
                lambda: nb_dataframe.train_dataframe(spark, data.train, data.vocab_size,
                                                     data.idx_to_label, num_partitions=parts)
            )
            acc_df, t_pred_df = _timeit(
                lambda: nb_dataframe.evaluate_dataframe(spark, model_df, data.test,
                                                        num_partitions=parts)
            )

            t_train_rdd_runs.append(t_train_rdd)
            t_pred_rdd_runs.append(t_pred_rdd)
            t_train_df_runs.append(t_train_df)
            t_pred_df_runs.append(t_pred_df)

            # Contrôle de cohérence : les deux versions donnent la même accuracy.
            assert abs(acc_rdd - acc_df) < 1e-9, (acc_rdd, acc_df)

        # Agrégation moyenne ± écart-type sur les ``repeat`` exécutions.
        m_train_rdd, s_train_rdd = _mean_std(t_train_rdd_runs)
        m_pred_rdd, s_pred_rdd = _mean_std(t_pred_rdd_runs)
        m_train_df, s_train_df = _mean_std(t_train_df_runs)
        m_pred_df, s_pred_df = _mean_std(t_pred_df_runs)

        common = dict(master=master, factor=factor, cores=sc.defaultParallelism,
                      repeat=repeat, n_train=n_train, n_test=n_test,
                      vocab_size=data.vocab_size)
        rows.append({**common, "version": "rdd",
                     "train_time_s": round(m_train_rdd, 4), "train_time_std": round(s_train_rdd, 4),
                     "predict_time_s": round(m_pred_rdd, 4), "predict_time_std": round(s_pred_rdd, 4),
                     "accuracy": round(acc_rdd, 4)})
        rows.append({**common, "version": "dataframe",
                     "train_time_s": round(m_train_df, 4), "train_time_std": round(s_train_df, 4),
                     "predict_time_s": round(m_pred_df, 4), "predict_time_std": round(s_pred_df, 4),
                     "accuracy": round(acc_df, 4)})

        print(f"  [{master} x{factor}] train={n_train} (repeat={repeat}) | "
              f"RDD train={m_train_rdd:.2f}±{s_train_rdd:.2f}s pred={m_pred_rdd:.2f}±{s_pred_rdd:.2f}s | "
              f"DF train={m_train_df:.2f}±{s_train_df:.2f}s pred={m_pred_df:.2f}±{s_pred_df:.2f}s | "
              f"acc={acc_rdd:.3f}")
    finally:
        spark.stop()

    return rows


def _plot(results: List[dict]) -> None:
    """Génère les graphiques PNG de scalabilité (avec barres d'erreur = écart-type)."""
    import matplotlib
    matplotlib.use("Agg")  # backend non interactif (pas de fenêtre)
    import matplotlib.pyplot as plt

    os.makedirs(RESULTS_DIR, exist_ok=True)

    # --- Graphe 1 : temps d'entraînement vs taille des données (à cores fixés) ---
    # On prend le master ayant le plus grand nombre de mesures par facteur.
    ref_master = max(
        {r["master"] for r in results},
        key=lambda m: len({r["factor"] for r in results if r["master"] == m}),
    )
    subset = [r for r in results if r["master"] == ref_master]
    if len({r["factor"] for r in subset}) >= 2:
        fig, ax = plt.subplots(figsize=(7, 5))
        for version in ("rdd", "dataframe"):
            pts = sorted((r for r in subset if r["version"] == version),
                         key=lambda r: r["n_train"])
            ax.errorbar([r["n_train"] for r in pts], [r["train_time_s"] for r in pts],
                        yerr=[r.get("train_time_std", 0.0) for r in pts],
                        marker="o", capsize=3, label=f"{version} (train)")
        ax.set_xlabel("Nombre de documents d'entraînement")
        ax.set_ylabel("Temps d'entraînement (s)")
        ax.set_title(f"Scalabilité en volume ({ref_master})")
        ax.legend()
        ax.grid(True, alpha=0.3)
        fig.tight_layout()
        fig.savefig(os.path.join(RESULTS_DIR, "scalability_datasize.png"), dpi=120)
        plt.close(fig)

    # --- Graphe 2 : temps vs nombre de cœurs (à facteur fixé le plus grand) ---
    factors_with_multi_cores = [
        f for f in {r["factor"] for r in results}
        if len({r["cores"] for r in results if r["factor"] == f}) >= 2
    ]
    if factors_with_multi_cores:
        ref_factor = max(factors_with_multi_cores)
        subset = [r for r in results if r["factor"] == ref_factor]
        fig, ax = plt.subplots(figsize=(7, 5))
        for version in ("rdd", "dataframe"):
            pts = sorted((r for r in subset if r["version"] == version),
                         key=lambda r: r["cores"])
            ax.errorbar([r["cores"] for r in pts], [r["train_time_s"] for r in pts],
                        yerr=[r.get("train_time_std", 0.0) for r in pts],
                        marker="o", capsize=3, label=f"{version} (train)")
        ax.set_xlabel("Nombre de cœurs (local[k])")
        ax.set_ylabel("Temps d'entraînement (s)")
        ax.set_title(f"Speed-up selon le parallélisme (facteur x{ref_factor})")
        ax.legend()
        ax.grid(True, alpha=0.3)
        fig.tight_layout()
        fig.savefig(os.path.join(RESULTS_DIR, "scalability_cores.png"), dpi=120)
        plt.close(fig)


def main() -> None:
    parser = argparse.ArgumentParser(description="Benchmark de scalabilité Naive Bayes/Spark")
    parser.add_argument("--factors", type=int, nargs="+", default=[1, 2, 4],
                        help="Facteurs de réplication du jeu de données.")
    parser.add_argument("--cores", type=int, nargs="+", default=[1, 2, 4],
                        help="Nombres de cœurs pour local[k].")
    parser.add_argument("--categories", type=str, nargs="+", default=None,
                        help="Sous-ensemble de catégories 20 Newsgroups (défaut: toutes).")
    parser.add_argument("--repeat", type=int, default=1,
                        help="Nombre d'exécutions par point ; les temps sont moyennés "
                             "(moyenne ± écart-type) pour lisser le bruit de mesure.")
    parser.add_argument("--quick", action="store_true",
                        help="Mode rapide : 4 catégories, facteurs [1,2], cores [1,2].")
    args = parser.parse_args()

    if args.repeat < 1:
        parser.error("--repeat doit être >= 1")

    if args.quick:
        args.categories = ["sci.space", "rec.autos", "comp.graphics", "talk.politics.misc"]
        args.factors = [1, 2]
        args.cores = [1, 2]

    print("Chargement de 20 Newsgroups ...")
    base_texts, base_labels = C.load_20newsgroups(categories=args.categories, subset="all")
    print(f"  {len(base_texts)} documents, {len(set(base_labels))} classes.")
    print(f"  repeat={args.repeat} exécution(s) par point.")

    all_rows: List[dict] = []

    # --- Expérience A : volume croissant à parallélisme maximal (local[*]) ---
    print("\n=== Scalabilité en VOLUME (local[*]) ===")
    for f in args.factors:
        all_rows += run_once("local[*]", f, base_texts, base_labels, repeat=args.repeat)

    # --- Expérience B : parallélisme croissant à volume fixe (facteur max) ---
    print("\n=== Scalabilité en PARALLÉLISME (facteur fixe) ===")
    fixed_factor = max(args.factors)
    for k in args.cores:
        all_rows += run_once(f"local[{k}]", fixed_factor, base_texts, base_labels,
                             repeat=args.repeat)

    # --- Export CSV -----------------------------------------------------------
    os.makedirs(RESULTS_DIR, exist_ok=True)
    csv_path = os.path.join(RESULTS_DIR, "results.csv")
    fieldnames = ["master", "cores", "factor", "repeat", "version", "n_train", "n_test",
                  "vocab_size", "train_time_s", "train_time_std",
                  "predict_time_s", "predict_time_std", "accuracy"]
    with open(csv_path, "w", newline="", encoding="utf-8") as fh:
        writer = csv.DictWriter(fh, fieldnames=fieldnames)
        writer.writeheader()
        writer.writerows(all_rows)
    print(f"\nMesures écrites dans {csv_path}")

    # --- Graphiques -----------------------------------------------------------
    _plot(all_rows)
    print(f"Graphiques écrits dans {RESULTS_DIR}/ (PNG).")


if __name__ == "__main__":
    main()
\end{lstlisting}
\subsection{\texttt{tests/test\_smoke.py}}
\begin{lstlisting}[style=py]
"""
test_smoke.py
=============

Smoke test : vérifie de bout en bout que les DEUX versions (RDD et DataFrames)
- s'entraînent sur un mini-jeu de 20 lignes,
- produisent EXACTEMENT le même modèle et donc la même accuracy.

Exécutable directement (``python tests/test_smoke.py``) ou via pytest
(``pytest tests/test_smoke.py``).
"""

from __future__ import annotations

import os
import sys

# Rendre le dossier src importable (nb_common, nb_rdd, nb_dataframe).
HERE = os.path.dirname(os.path.abspath(__file__))
SRC = os.path.join(os.path.dirname(HERE), "src")
sys.path.insert(0, SRC)

import nb_common as C          # noqa: E402
import nb_rdd                   # noqa: E402
import nb_dataframe            # noqa: E402


# --- Mini-jeu de données : 20 messages étiquetés spam/ham -------------------
# Volontairement séparable pour que la prédiction ait du sens sur si peu de data.
MINI_TEXTS = [
    "win a free prize now",
    "free money win cash prize",
    "claim your free lottery prize",
    "congratulations you won free cash",
    "urgent free offer claim now",
    "win big money free entry",
    "free tickets claim your prize",
    "cash prize winner claim now",
    "exclusive free offer win now",
    "you won a free vacation prize",
    "hey are we still meeting today",
    "call me when you get home",
    "lunch tomorrow at noon sounds good",
    "can you send me the report please",
    "happy birthday see you tonight",
    "the meeting is moved to friday",
    "thanks for your help yesterday",
    "i will be late for dinner",
    "did you finish the homework yet",
    "see you at the office monday",
]
MINI_LABELS = (["spam"] * 10) + (["ham"] * 10)


def run() -> None:
    # Démarrage d'une SparkSession locale utilisant tous les cœurs (local[*]).
    spark = C.get_spark(app_name="SmokeTest", master="local[*]")
    spark.sparkContext.setLogLevel("ERROR")  # logs Spark discrets
    sc = spark.sparkContext

    try:
        # Découpage train/test reproductible, puis vectorisation commune.
        X_tr, X_te, y_tr, y_te = C.train_test_split_texts(
            MINI_TEXTS, MINI_LABELS, test_size=0.3, seed=0
        )
        data = C.prepare(X_tr, y_tr, X_te, y_te)

        # Entraînement des deux versions sur EXACTEMENT les mêmes données.
        model_rdd = nb_rdd.train_rdd(
            sc, data.train, data.vocab_size, data.idx_to_label
        )
        model_df = nb_dataframe.train_dataframe(
            spark, data.train, data.vocab_size, data.idx_to_label
        )

        # Évaluation.
        acc_rdd = nb_rdd.evaluate_rdd(sc, model_rdd, data.test)
        acc_df = nb_dataframe.evaluate_dataframe(spark, model_df, data.test)

        # Prédictions brutes (doivent être identiques ligne à ligne).
        pred_rdd = nb_rdd.predict_rdd(sc, model_rdd, data.test)
        pred_df = nb_dataframe.predict_dataframe(spark, model_df, data.test)

        print(f"Taille vocabulaire : {data.vocab_size}")
        print(f"Train / Test       : {len(data.train)} / {len(data.test)}")
        print(f"Accuracy RDD       : {acc_rdd:.4f}")
        print(f"Accuracy DataFrame : {acc_df:.4f}")
        print(f"Prédictions RDD    : {pred_rdd}")
        print(f"Prédictions DF     : {pred_df}")

        # --- Assertions du smoke test ---------------------------------------
        assert acc_rdd == acc_df, (
            f"Accuracies différentes : RDD={acc_rdd} vs DF={acc_df}"
        )
        assert pred_rdd == pred_df, (
            "Les prédictions ligne à ligne diffèrent entre RDD et DataFrame."
        )
        print("\nOK : RDD et DataFrame donnent des prédictions et une accuracy identiques.")
    finally:
        spark.stop()


def test_smoke():
    """Point d'entrée pytest."""
    run()


if __name__ == "__main__":
    run()
\end{lstlisting}
\subsection{\texttt{data/download\_sms\_spam.py}}
\begin{lstlisting}[style=py]
"""
download_sms_spam.py
====================

Télécharge le jeu de données **SMS Spam Collection** (UCI) et l'exporte en CSV
(``data/sms_spam.csv``) avec deux colonnes : ``label`` (ham/spam) et ``text``.

C'est le PETIT jeu de données, utilisé par le notebook de démonstration et le
smoke test. Il ne nécessite pas de cluster.

Usage :
    python data/download_sms_spam.py
"""

from __future__ import annotations

import csv
import io
import os
import urllib.request
import zipfile

# Miroir officiel UCI de la SMS Spam Collection (fichier zip contenant un TSV).
URL = "https://archive.ics.uci.edu/static/public/228/sms+spam+collection.zip"
# URL de repli (autre miroir historique).
FALLBACK_URL = (
    "https://raw.githubusercontent.com/justmarkham/pycon-2016-tutorial/master/"
    "data/sms.tsv"
)

HERE = os.path.dirname(os.path.abspath(__file__))
OUT_CSV = os.path.join(HERE, "sms_spam.csv")


def _rows_from_zip(raw: bytes):
    """Extrait les lignes (label, text) du zip UCI (fichier 'SMSSpamCollection')."""
    with zipfile.ZipFile(io.BytesIO(raw)) as zf:
        # Le zip contient un fichier tabulé : "<label>\t<message>" par ligne.
        name = next(n for n in zf.namelist() if "SMSSpamCollection" in n)
        with zf.open(name) as fh:
            for line in io.TextIOWrapper(fh, encoding="utf-8"):
                line = line.rstrip("\n")
                if not line:
                    continue
                label, _, text = line.partition("\t")
                yield label.strip(), text


def _rows_from_tsv(raw: bytes):
    """Extrait les lignes (label, text) du miroir TSV de repli."""
    for line in io.StringIO(raw.decode("utf-8")):
        line = line.rstrip("\n")
        if not line:
            continue
        label, _, text = line.partition("\t")
        yield label.strip(), text


def main() -> None:
    print(f"Téléchargement depuis {URL} ...")
    try:
        with urllib.request.urlopen(URL, timeout=60) as resp:
            raw = resp.read()
        rows = list(_rows_from_zip(raw))
    except Exception as exc:  # pragma: no cover - dépend du réseau
        print(f"  échec ({exc}); tentative sur le miroir de repli...")
        with urllib.request.urlopen(FALLBACK_URL, timeout=60) as resp:
            raw = resp.read()
        rows = list(_rows_from_tsv(raw))

    # Écriture du CSV normalisé (en-tête label,text).
    with open(OUT_CSV, "w", encoding="utf-8", newline="") as fh:
        writer = csv.writer(fh)
        writer.writerow(["label", "text"])
        writer.writerows(rows)

    n_spam = sum(1 for lab, _ in rows if lab == "spam")
    print(f"OK : {len(rows)} messages écrits dans {OUT_CSV} "
          f"({n_spam} spam / {len(rows) - n_spam} ham).")


if __name__ == "__main__":
    main()
\end{lstlisting}
\subsection{\texttt{data/download\_20newsgroups.py}}
\begin{lstlisting}[style=py]
"""
download_20newsgroups.py
========================

Pré-télécharge (met en cache) le jeu de données **20 Newsgroups** via
scikit-learn. C'est le GRAND jeu de données utilisé pour l'étude de scalabilité
(``src/benchmark.py``).

scikit-learn met les données en cache dans ``~/scikit_learn_data`` ; ce script
force simplement ce téléchargement une fois pour toutes et affiche un résumé.

Usage :
    python data/download_20newsgroups.py
"""

from __future__ import annotations

from sklearn.datasets import fetch_20newsgroups


def main() -> None:
    print("Téléchargement / mise en cache de 20 Newsgroups via scikit-learn ...")
    # subset="all" = train + test réunis. remove=(...) enlève en-têtes/pieds/citations
    # pour ne garder qu'un signal textuel (évite un sur-apprentissage trivial).
    bunch = fetch_20newsgroups(
        subset="all",
        remove=("headers", "footers", "quotes"),
        shuffle=True,
        random_state=42,
    )
    print(f"OK : {len(bunch.data)} documents, {len(bunch.target_names)} catégories.")
    print("Catégories :")
    for name in bunch.target_names:
        print(f"  - {name}")
    print("\nLes données sont en cache (~/scikit_learn_data) et prêtes pour le benchmark.")


if __name__ == "__main__":
    main()
\end{lstlisting}
